\chapter{Related Works}
\label{ch:related}

Work relevant to this proposal falls into four families: (1)
generate-and-validate repair and the size of its validation bill, (2)
regression test-case prioritisation, which is where the scheduling machinery
comes from, (3) prioritisation specialised to repair loops, which is where the
direct baseline lives, and (4) LLM-based repair loops and their token
economics, which is where the proposer-side half of the design belongs. This
chapter walks the four, consolidates them into one table, and distils the gaps.

\section{Survey of Representative Works}
\label{sec:survey}

\paragraph{Generate-and-validate repair and the validation bill.}
\textsc{GenProg}~\citep{legoues2012genprog} established the loop that the rest
of this literature refines: mutate, validate against the test suite, keep what
survives. Its systematic study repaired $55$ of $105$ real
bugs~\citep{legoues2012systematic}, and validation dominates its runtime,
because each candidate is run against the suite from scratch. The structural weakness of validating against a
test suite was made precise by \citet{qi2015plausibility}, who re-examined the
patches reported by three generate-and-validate systems and found that most
were \emph{plausible} rather than correct --- \textsc{GenProg} produced
correct patches for $2$ defects out of $18$ for which it produced a plausible
patch, \textsc{AE} $3$ of $27$, and \textsc{RSRepair} $2$ of $10$. That result
is the reason this proposal treats the strength of the oracle (how many probes
a patch must survive) as a parameter to be stated rather than an implementation
detail, and the reason the soundness invariant of
Section~\ref{sec:soundness} is asserted on every cell.
Defects4J~\citep{just2014defects4j} supplied the benchmark the family
standardised on; ConDefects~\citep{wu2024condefects} supplies the one used
here, for the data-leakage reason given in Section~\ref{sec:scope}.

\paragraph{Regression test-case prioritisation.}
\citet{rothermel2001prioritizing} defined the problem and the two greedy
families that still dominate it: \emph{total} strategies, which rank a test by
how much it covers, and \emph{additional} strategies, which rank by how much it
covers that nothing scheduled so far has covered. The
survey of \citet{lou2019survey} organises two decades of refinements of those
two. Crucially for this proposal, \citet{elbaum2001cost} observed early that
the field's standard metric assumed uniform test costs and uniform fault
severities, and proposed a cost-cognizant replacement. That paper is the
philosophical ancestor of the $\hat{m}$ term here; what it did not have, and
what this work adds, is an \emph{online} cost estimate learned inside the loop
rather than a cost supplied in advance.

\paragraph{Prioritisation inside a repair loop.}
\citet{qi2013frtp} introduced fault-recorded testing prioritisation (\frtp),
implemented as \textsc{TrpAutoRepair}, and evaluated it against
\textsc{GenProg} on five subject programs with sixteen real bugs: it matched or
beat \textsc{GenProg}'s success rate and improved repair efficiency in fifteen
of the sixteen. \frtp\ ranks a test case by its \emph{failing count} --- the
number of validated patches the case has already falsified --- and gathers that
information from the repair process itself rather than from a separate
profiling pass. It is the published baseline for this work, and the arm named
\code{fault-record} in Chapter~\ref{ch:evaluation} implements exactly that
ranking, not a weakened variant. \citet{venugopal2020mpp} extend the idea with
modification-point awareness, maintaining a table whose entries are the
``patch killing count'' of a test case against a program variant, which is the
same quantity under another name.

The most thorough measurement of patch-validation cost is
\citet{lou2024aprregression}, who study $2{,}532{,}915$ patches produced by
twelve APR systems on Defects4J. Two of their findings bear directly on this
proposal. First, they report results in test executions rather than seconds,
stating that executions are ``more stable and more suitable for future work to
reproduce''. Second, they find that \emph{selection} dominates
\emph{prioritisation} in their setting: class-level regression test selection
alone reduces executions by about $37\,\%$ for \textsc{PraPR}, where the best
prioritisation approach reduces them by only about $15\,\%$. Both findings are
properties of their setting rather than of the problem, and
Section~\ref{sec:gap} explains why neither transfers to the one studied here.

\paragraph{LLM repair loops and their token economics.}
\textsc{ChatRepair}~\citep{xia2024chatrepair} made the conversational loop
standard, repairing $162$ of $337$ Defects4J bugs at roughly \$0.42 each by
feeding failing-test feedback back into the conversation.
\textsc{CigaR}~\citep{hidvegi2024cigar} attacked its cost directly and reports
$127$K tokens per bug against a reimplemented \textsc{ChatRepair} baseline's
$467$K, a $73\,\%$ reduction, while producing correct patches for $171$ of
$429$ bugs. The sharpest result in the family is negative:
\citet{chen2025conversational} hold the budget fixed at $24$ model calls per
bug and find that \emph{independent repeated prompting} repairs $23$ bugs where
\textsc{ChatRepair} repairs $18$, and that $65\,\%$ of \textsc{ChatRepair}'s
generated patches are duplicates. The obvious iterative strategy is therefore
already published as worse than simply drawing again, which fixes the baseline
any steering arm in this project must beat.

\paragraph{The immediate predecessor.}
\cegmem\ --- the author's previous study, described in detail in
Section~\ref{sec:cegmem-premise} --- sits between the two halves. It keyed a
memory of refuted patches by a failure type $\theta = (\text{edit location})
\times (\text{violated property})$ and used it both to block redundant
candidates before the oracle ran and to steer the proposer away from exhausted
classes. On $309$ ConDefects tasks it cut oracle calls per episode by
$7.07\times$ and raised the repair rate from $0.526$ to $0.581$. It is cited
here as the premise of this proposal rather than as a competitor.

\section{Comparison Table}
\label{sec:comparison}

Table~\ref{tab:related} consolidates the surveyed systems against the four
design dimensions this proposal cares about: \textbf{M1} what the memory is
keyed by, \textbf{M2} what quantity is ranked, \textbf{M3} whether the cost of
a probe enters that quantity, and \textbf{M4} whether the evidence is returned
to the proposer.

\begin{table}[H]
\centering
\caption{Representative works against the four dimensions. ``---'' means the
dimension is absent from the system.}
\label{tab:related}
\scriptsize
\setlength{\tabcolsep}{4pt}
\begin{tabularx}{\textwidth}{@{}r l l l l c c X@{}}
\toprule
\# & System & Yr & M1 key & M2 ranked quantity & M3 & M4 & Headline result\\
\midrule
1 & \textsc{GenProg}~\citep{legoues2012genprog,legoues2012systematic} & 2012 & ---  & --- (suite order)      &    &    & 55/105 bugs at \$8 each\\
2 & \frtp~\citep{qi2013frtp}                           & 2013 & failure class  & kill \emph{count}      &    &    & Faster than GenProg on 15/16\\
3 & Cost-cognizant TCP~\citep{elbaum2001cost}          & 2001 & ---            & coverage / cost        & \checkmark &  & $\mathrm{APFD}_c$ metric\\
4 & Additional greedy~\citep{rothermel2001prioritizing}& 2001 & ---            & marginal coverage      &    &    & The $\mathrm{RTP_{add}}$ family\\
5 & MPP~\citep{venugopal2020mpp}                       & 2020 & modif.\ point  & patch killing count    &    &    & Faster patch validation\\
6 & APR\,$\times$\,RTS~\citep{lou2024aprregression}    & 2024 & ---            & selection + priority   &    &    & 2.53M patches; sel.\ $>$ prior.\\
7 & \textsc{ChatRepair}~\citep{xia2024chatrepair}      & 2024 & transcript     & --- (suite order)      &    & \checkmark & 162/337 at \$0.42/bug\\
8 & \textsc{CigaR}~\citep{hidvegi2024cigar}            & 2024 & transcript     & --- (suite order)      &    & \checkmark & 127K vs 467K tokens/bug\\
9 & Repeated sampling~\citep{chen2025conversational}   & 2025 & none           & ---                    &    &    & 23 vs 18 bugs; 65\% dupes\\
10& \cegmem\ (premise)                                 & 2026 & failure class  & --- (pool order)       &    & \checkmark & $7.07\times$ fewer oracle calls\\
\midrule
-- & \textbf{\caster\ (this work)}                     & 2026 & \textbf{test case} & \textbf{rate / cost} & \checkmark & \checkmark & See Sec.~\ref{sec:pilot}\\
\bottomrule
\end{tabularx}
\end{table}

\section{Research Gap}
\label{sec:gap}

Three gaps follow from the table, and the third is the one the pilot measures.

\paragraph{Gap 1 --- the memory is keyed by a class, so it cannot schedule.}
Every memory in rows 2, 5, 7, 8 and 10 is indexed by something that is not a
test case: a failure class, a modification point, or a conversation transcript.
Such an index answers ``has this kind of patch been refuted before?''. It
cannot answer ``which probe should run first?'', because the probe is not the
thing being indexed. \cegmem\ is the clean demonstration: its class index
removed $86\,\%$ of oracle calls and still left every surviving call walking
the pool in alphabetical order. Re-keying by test case is what makes one record
serve the scheduler and the prompt at once.

\paragraph{Gap 2 --- the ranked quantity discards the evidence of \emph{not}
killing.} \frtp\ and its descendants rank on a count of kills. A count cannot
distinguish a case that has failed to refute ten patches from a case that has
never been run, because both have count zero, and the distinction is exactly
what tells you whether to keep paying for the case. A \emph{rate} --- kills
over attempts --- is the first quantity that reads the denominator, and the
first that can decay a case whose early record was luck.

\paragraph{Gap 3 --- the cost of a probe is not in the score, and on
heterogeneous pools that is the term that matters.} Only
\citet{elbaum2001cost} put cost in the score at all, and they assume it is
known in advance rather than learned online. This gap is usually invisible
because of the unit of account: \citet{lou2024aprregression} measure
executions, for the reproducibility reason quoted above, and under that unit a
cheap probe and an expensive one are the same event. The measured cost
distribution of Section~\ref{sec:pilot-arena} shows what that hides here: a
$58.7\times$ spread between the median and the 99th percentile, \emph{inside}
the same pool. The pilot's headline (Section~\ref{sec:pilot-measured}) is that
the two units disagree in sign --- \km\ pays $23.3\,\%$ more executions than
\frtp\ for $15.7\,\%$ less CPU --- so a study that reports executions alone
would conclude \frtp\ wins.

Two properties of the setting are worth stating, because they explain why
\citeauthor{lou2024aprregression}'s ``selection beats prioritisation'' finding
does not transfer. First, their subjects are Java programs with method-level
test suites, where regression test selection can skip a test because no changed
method reaches it. Here a probe is a \emph{whole-program execution on one
input}: there is no call graph to cut and nothing to select away, so
prioritisation is the only lever available. Second, their cost unit is uniform
by construction; here it is not.

\section{Direct Competitors}
\label{sec:competitors}

\begin{table}[H]
\centering
\caption{Direct competitors and the differentiator in each case.}
\label{tab:competitors}
\small
\begin{tabularx}{\textwidth}{@{}l >{\raggedright\arraybackslash}p{0.17\textwidth} X X@{}}
\toprule
Rank & Competitor & Similarity & Our differentiator\\
\midrule
\#1 & \frtp~\citep{qi2013frtp} & Same place in the loop; a memory that reorders probes & A rate over attempts, not a count of kills; and divided by a measured cost. Pilot: $-15.7\,\%$ CPU against it\\
\#2 & \cegmem\ (premise) & Same corpus, same oracle, same proposer & Keyed by test case, not failure class; schedules the probe rather than blocking the patch\\
\#3 & MPP~\citep{venugopal2020mpp} & ``Patch killing count'' table & Same gap as \frtp: a count, no denominator, no cost\\
\#4 & Cost-cognizant TCP~\citep{elbaum2001cost} & Cost appears in the score & Cost is \emph{learned online} from \code{getrusage} inside the loop, not supplied in advance\\
\#5 & Additional greedy~\citep{rothermel2001prioritizing} & A scheduling policy over the same pool & Measured here as an arm (Sec.~\ref{sec:pilot-replay}) and found no better than pool order on this corpus\\
\#6 & APR\,$\times$\,RTS~\citep{lou2024aprregression} & Patch-validation cost, at scale & Selection is unavailable in this setting; and we report CPU, where they report executions\\
\#7 & \textsc{CigaR}~\citep{hidvegi2024cigar} & Explicit cost-efficiency goal & Their currency is tokens, ours is oracle CPU; the two are complementary halves of the same bill\\
\#8 & Repeated sampling~\citep{chen2025conversational} & The strong baseline on the proposer side & It is our RQ3 baseline, not our competitor on RQ1; scheduling is orthogonal to how the draw is made\\
\bottomrule
\end{tabularx}
\end{table}
